\section{Расчётная задача 7}

Пусть~$V$ --- четырёхмерное векторное пространство над~$\RR$, 
\[
	\hh = (h_1, h_2, h_3, h_4) \subset V
\]
--- ортонормированный базис. Пусть также пространство~$V$ евклидово со стандартным скалярным произведением:~$G = \diag(1, 1, 1, 1)$, то есть
\[
	g_{ij} = \spr{e_i}{e_j} = \delta_{ij}
	.
\]

Провести процесс ортогонализации Грама--Шмидта с заданной системой векторов~$\uu = (u_1, u_2, u_3, u_4)$, где
\begin{align*}
	u_1 &=
	\hh
	\Vect{-1 \\ 5 \\ 2 \\ -1}
	,
	&
	u_2 &=
	\hh
	\Vect{3 \\ 7 \\ -2 \\ -3}
	,
	&
	u_3 &=
	\hh
	\Vect{2 \\ 4 \\ 4 \\ -5}
	,
	&
	u_4 &=
	\hh
	\Vect{10 \\ 7 \\ 9 \\ 12}
	.
\end{align*}

\begin{solution}
	Организуем координаты векторов~$\uu$ в матрицу~$U$, так что $\uu = \hh U$:
	\[
		U =
		\left( 
			\begin{array}{r|r|r|r}
				-1 & 3 & 2 & 10
				\\
				5 & 7 & 4 & 7
				\\
				2 & -2 & 4 & 9
				\\
				-1 & -3 & -5 & 12
			\end{array}
		\right)
		.
	\]
	Как водится, будем обозначать символом~$g_\nu$ вектор, ортогональный предыдущим~$g_\mu$ при~$\mu < \nu$; а~$e_\nu := g_\nu / \norm{g_\nu}$; при этом $\norm{v}^2 := \spr{v}{v}$.

	Формула для расчёта~$g_\nu$ имеет вид
	\[
		g_\nu =
		u_\nu - \pr_{L_{\nu-1}}(u_\nu),
	\]
	где~$L_{\nu-1} = \opspan_\RR(g_1, \dots, g_{\nu-1})$,
	\[
		\pr_{L_{\nu-1}}(u_\nu) =
		\sum_{i = 1}^{\nu-1}
			\pr_{L_i}(u_\nu)
		=
		\sum_{i = 1}^{\nu-1}
			\frac{\spr{u_\nu}{g_i}}{\spr{g_i}{g_i}}\,
			g_i
		;
	\]
	так что
	\begin{equation}
		\label{eq:num_07:g}
		g_\nu =
		u_\nu -
		\sum_{i=1}^{\nu-1}
			\frac{\spr{u_\nu}{g_i}}{\spr{g_i}{g_i}}\,
			g_i
		.
	\end{equation}

	\paragraph{Первый вектор.}
	Положим~$g_1 := u_1$. Тогда
	\[
		\norm{g_1}^2 =
		(-1)^2 + 5^2 + 2^2 + (-1)^2 =
		31.
	\]
	Поэтому
	\[
		e_1 =
		\hh
		\cdot
		\frac{1}{\sqrt{31}}
		\Vect{-1 \\ 5 \\ 2 \\ -1}
		.
	\]

	\paragraph{Второй вектор.}
	Найдём~$g_2$ по формуле~\eqref{eq:num_07:g}:
	\[
		g_2 =
		u_2 -
		\vproj{u_2}{g_1}
		.
	\]
	Найдём~$\spr{u_2}{g_1}$:
	\[
		\spr{u_2}{g_1} =
		\spr{
			\hh
			\Vect{3 \\ 7 \\ -2 \\ -3}
		}{
			\hh
			\Vect{-1 \\ 5 \\ 2 \\ -1}
		}
		=
		-3 + 35 - 4 + 3
		=
		31
		.
	\]
	Тогда
	\[
		g_2 =
		u_2 -
		\frac{31}{31}\,g_1
		=
		u_2 - g_1
		=
		\hh
		\left[ 
			\Vect{3 \\ 7 \\ -2 \\ -3}
			-
			\Vect{-1 \\ 5 \\ 2 \\ -1}
		\right]
		=
		\hh
		\Vect{4 \\ 2 \\ -4 \\ -2}
		=
		2
		\hh
		\Vect{2 \\ 1 \\ -2 \\ -1}
		.
	\]
	Найдём квадрат нормы~$g_2$:
	\[
		\spr{g_2}{g_2} =
		\spr{
			2
			\hh
			\Vect{2 \\ 1 \\ -2 \\ -1}
		}{
			2
			\hh
			\Vect{2 \\ 1 \\ -2 \\ -1}
		}
		=
		4 \cdot (4 + 1 + 4 + 1)
		=
		4 \cdot 10 = 40
		.
	\]
	Тогда
	\[
		e_2 =
		\hh
		\cdot
		\frac{1}{2\sqrt{10}}
		\cdot
		2
		\Vect{2 \\ 1 \\ -2 \\ -1}
		=
		\hh
		\cdot
		\frac{1}{\sqrt{10}}
		\Vect{2 \\ 1 \\ -2 \\ -1}
		.
	\]

	\paragraph{Третий вектор.}
	Формула для нахождения~$g_3$ имеет вид
	\[
		g_3 =
		u_3 -
		\vproj{u_3}{g_1}
		-
		\vproj{u_3}{g_2}
		.
	\]
	Поэтому найдём скалярные произведения~$\spr{u_3}{g_1}$ и~$\spr{u_3}{g_2}$:
	\begin{align*}
		\spr{u_3}{g_1} &=
		\spr{
			\hh
			\Vect{2 \\ 4 \\ 4 \\ -5}
		}{
			\hh
			\Vect{-1 \\ 5 \\ 2 \\ -1}
		}
		=
		-2 + 20 + 8 + 5 = 31
		;
		\\
		\spr{u_3}{g_2} &=
		\spr{
			\hh
			\Vect{2 \\ 4 \\ 4 \\ -5}
		}{
			2
			\hh
			\Vect{2 \\ 1 \\ -2 \\ -1}
		}
		=
		2 \cdot (4 + 4 - 8 + 5)
		=
		10
		.
	\end{align*}
	Тогда
	\begin{multline*}
		g_3 =
		\hh
		\Vect{2 \\ 4 \\ 4 \\ -5}
		-
		\frac{31}{31}
		\hh
		\Vect{-1 \\ 5 \\ 2 \\ -1}
		-
		\frac{10}{40}
		\cdot
		2
		\hh
		\Vect{2 \\ 1 \\ -2 \\ -1}
		=
		\hh
		\left[ 
			\Vect{2 \\ 4 \\ 4 \\ -5}
			-
			\Vect{-1 \\ 5 \\ 2 \\ -1}
			-
			\frac{1}{2}
			\Vect{2 \\ 1 \\ -2 \\ -1}
		\right]
		\nline{=}
		\hh
		\cdot
		\frac{1}{2}
		\left[ 
			\Vect{4 \\ 8 \\ 8 \\ -10}
			-
			\Vect{-2 \\ 10 \\ 4 \\ -2}
			-
			\Vect{2 \\ 1 \\ -2 \\ -1}
		\right]
		=
		\hh
		\cdot
		\frac{1}{2}
		\Vect{4 \\ -3 \\ 6 \\ -7}
		.
	\end{multline*}
	Тогда
	\[
		\spr{g_3}{g_3} =
		\frac{1}{4}
		\cdot
		(16 + 9 + 36 + 49)
		=
		% \frac{(3^2 + 4^2 + 5^2 + 6^2 + 7^2) - 5^2}{2^2}
		% =
		% \frac{(1^2 + \dots + 7^2) - (1^2 + 2^2 + 5^2)}{2^2}
		% =
		% \frac{\bigl( 7 \cdot (7 + 1/2) \cdot 8 \bigr)/3 - (1^2 + 2^2 + 5^2)}{2^2}
		% =
		% \frac{\bigl( 7 \cdot (14 + 1) \cdot 8 \bigr)/6 - (1^2 + 2^2 + 5^2)}{2^2}
		% =
		% \frac{\bigl( 7 \cdot 15 \cdot 8 \bigr)/6 - (1^2 + 2^2 + 5^2)}{2^2}
		\frac{55}{2}
		.
	\]

	\paragraph{Четвёртый вектор.}
	Вектор~$g_4$ находим по формуле
	\[
		g_4 =
		u_4 -
		\vproj{u_4}{g_1} -
		\vproj{u_4}{g_2} -
		\vproj{u_4}{g_3}
		.
	\]
	Найдём скалярные произведения~$\spr{u_4}{g_\mu}$, $\mu \in \set*{1, 2, 3}$:
	\begin{align*}
		\spr{u_4}{g_1} &=
		\spr{
			\hh
			\Vect{10 \\ 7 \\ 9 \\ 12}
		}{
			\hh
			\Vect{-1 \\ 5 \\ 2 \\ -1}
		}
		=
		-10 + 35 + 18 - 12
		=
		31
		;
		\\
		\spr{u_4}{g_2} &=
		\spr{
			\hh
			\Vect{10 \\ 7 \\ 9 \\ 12}
		}{
			2
			\hh
			\Vect{2 \\ 1 \\ -2 \\ -1}
		}
		=
		2 \cdot (20 + 7 - 18 - 12)
		=
		-6
		;
		\\
		\spr{u_4}{g_3} &=
		\spr{
			\hh
			\Vect{10 \\ 7 \\ 9 \\ 12}
		}{
			\hh
			\cdot
			\frac{1}{2}
			\Vect{4 \\ -3 \\ 6 \\ -7}
		}
		=
		\frac{1}{2}
		\cdot
		\bigl(
			40 + 7 \cdot (-3) + 9 \cdot 6 + 12 \cdot (-7)
		\bigr)
		\nlinea{=}
		\frac{1}{2}
		\cdot
		\bigl(
			40 - 21 + 54 - 84
		\bigr)
		=
		-\frac{11}{2}
		.
	\end{align*}
	Тогда
	\begin{multline*}
		g_4
		=
		\hh
		\Vect{10 \\ 7 \\ 9 \\ 12}
		-
		\frac{31}{31}
		\cdot
		\hh
		\Vect{-1 \\ 5 \\ 2 \\ -1}
		-
		\frac{-6}{40}
		\cdot
		2
		\cdot
		\hh
		\Vect{2 \\ 1 \\ -2 \\ -1}
		-
		\frac{-11/2}{55/2}
		\cdot
		\frac{1}{2}
		\cdot
		\hh
		\Vect{4 \\ -3 \\ 6 \\ -7}
		\nline{=}
		\hh
		\left[ 
			\Vect{10 \\ 7 \\ 9 \\ 12}
			-
			\Vect{-1 \\ 5 \\ 2 \\ -1}
			+
			\frac{3}{10}
			\Vect{2 \\ 1 \\ -2 \\ -1}
			+
			\frac{1}{10}
			\Vect{4 \\ -3 \\ 6 \\ -7}
		\right]
		\nline{=}
		\hh
		\cdot
		\frac{1}{10}
		\left[ 
			10
			\Vect{10 \\ 7 \\ 9 \\ 12}
			-
			10
			\Vect{-1 \\ 5 \\ 2 \\ -1}
			+
			3
			\Vect{2 \\ 1 \\ -2 \\ -1}
			+
			\Vect{4 \\ -3 \\ 6 \\ -7}
		\right]
		\nline{=}
		\hh
		\cdot
		\frac{1}{10}
		\left[ 
			\Vect{100 \\ 70 \\ 90 \\ 120}
			-
			\Vect{-10 \\ 50 \\ 20 \\ -10}
			+
			\Vect{6 \\ 3 \\ -6 \\ -3}
			+
			\Vect{4 \\ -3 \\ 6 \\ -7}
		\right]
		=
		\hh
		\cdot
		\frac{1}{10}
		\Vect{
			120 \\
			20 \\
			70 \\
			120
		}
		=
		\hh
		\cdot
		\Vect{
			12 \\ 2 \\ 7 \\ 12
		}
		.
	\end{multline*}
	Теперь вычислим квадрат нормы:
	\[
		\spr{g_2}{g_2} =
		12^2 + 2^2 + 7^2 + 12^2 =
		2 \cdot 144 + 4 + 49 =
		288 + 50 + 3 =
		290 + 50 + 1 =
		341
		.
	\]
	Поэтому
	\[
		e_4 =
		\frac{1}{\sqrt{341}}
		\cdot
		\hh
		\Vect{
			12 \\ 2 \\ 7 \\ 12
		}
		.
	\]
	Поэтому
	\[
		\bgg = \hh \cdot G =
		\hh \cdot
		\left( 
			\begin{array}{r|r|r|r}
				-1 &  4 &  2   & 12 \\
				 5 &  2 & -3/2 &  2 \\
				 2 & -4 &  3   &  7 \\
				-1 & -2 & -7/2 & 12
			\end{array}
		\right)
		=
		\hh \cdot
		\left( 
			\begin{array}{r|r|r|r}
				-1 &  2 &  4 & 12 \\
				 5 &  1 & -3 &  2 \\
				 2 & -2 &  6 &  7 \\
				-1 & -1 & -7 & 12
			\end{array}
		\right)
		\begin{pmatrix}
			1 \\
			& 2 \\
			&& 1/2 \\
			&&& 1
		\end{pmatrix}
		.
	\]
	(Умножение на диагональную матрицу $\diag(\lambda_1, \dots, \lambda_n)$ справа умножает на~$\lambda_i$ $i$-ый столбец.)

	Поэтому
	\[
		\ee = \hh \cdot G \cdot N,
	\]
	где $N = \diag(1/\norm{g_1}, \dots, 1/\norm{g_4})$:
	\begin{multline*}
		\hh \cdot
		\left( 
			\begin{array}{r|r|r|r}
				-1 &  2 &  4 & 12 \\
				 5 &  1 & -3 &  2 \\
				 2 & -2 &  6 &  7 \\
				-1 & -1 & -7 & 12
			\end{array}
		\right)
		\begin{pmatrix}
			1 \\
			& 2 \\
			&& 1/2 \\
			&&& 1
		\end{pmatrix}
		\begin{pmatrix}
			1/\sqrt{31} \\
			& 2\sqrt{10} \\
			&& \sqrt{55/2} \\
			&&& \sqrt{341}
		\end{pmatrix}
		\nline{=}
		\hh \cdot
		\left( 
			\begin{array}{r|r|r|r}
				-1 &  2 &  4 & 12 \\
				 5 &  1 & -3 &  2 \\
				 2 & -2 &  6 &  7 \\
				-1 & -1 & -7 & 12
			\end{array}
		\right)
		\begin{pmatrix}
			1/\sqrt{31} \\
			& 4\sqrt{10} \\
			&& \sqrt{55/8} \\
			&&& \sqrt{341}
		\end{pmatrix}
		.
		% \proofmark
	\end{multline*}
\end{solution}
